\begin{frame}{Fundamento de los ratios}
\small
\begin{block}{Definición}
Un ratio financiero es una relación matemática entre partidas de los estados financieros que permite evaluar una dimensión específica de la situación o del desempeño empresarial.
\end{block}

\begin{columns}[T,totalwidth=\textwidth]
\column{0.52\textwidth}
\begin{tabular}{ll}
\toprule
Categoría & Qué evalúa \\
\midrule
Liquidez & Capacidad de pago de corto plazo \\
Gestión & Eficiencia en el uso de recursos \\
Solvencia & Nivel de deuda y financiamiento \\
Rentabilidad & Capacidad para generar utilidad \\
\bottomrule
\end{tabular}

\column{0.43\textwidth}
\begin{alertblock}{Advertencia académica}
Un ratio no debe calificarse automáticamente como favorable o desfavorable; debe compararse con periodos anteriores, políticas internas y características del sector.
\end{alertblock}

\begin{center}
\begin{tikzpicture}[node distance=0.25cm, every node/.style={font=\scriptsize}]
\node[draw, rounded corners=2pt] (a) {Datos};
\node[draw, rounded corners=2pt, right=of a] (b) {Cálculo};
\node[draw, rounded corners=2pt, below=of b] (c) {Resultado};
\node[draw, rounded corners=2pt, left=of c] (d) {Interpretación};
\node[draw, rounded corners=2pt, below=of d] (e) {Decisión};
\draw[-{Stealth}] (a) -- (b);
\draw[-{Stealth}] (b) -- (c);
\draw[-{Stealth}] (c) -- (d);
\draw[-{Stealth}] (d) -- (e);
\end{tikzpicture}
\end{center}
\end{columns}
\end{frame}
